\documentclass[../../../main]{subfiles}

\ifSubfilesClassLoaded{%
    \zexternaldocument*{../../1.Category-Theory/exported/aux/main}%
    \zexternaldocument*{../../2.Homological-Algebra/exported/aux/main}%
    \zexternaldocument*{../../3.Sheaves/1.Sheaves/exported/aux/main}%
}{}

\begin{document}
\section{Sheaves of modules} \label{sec:sheaves-of-modules}

\begin{definition} \label{def:presheaf-of-modules}
  Let \(X\) be a topological space.
  Let \(\mathcal{O}_{X}\) be a \hyperref[def:presheaf]{presheaf} of rings over \(X\).

  A \textbf{presheaf of \(\mathcal{O}_{X}\)-modules} is a \hyperref[def:presheaf]{presheaf} of abelian groups \(\mathcal{F}\) together with a map of \hyperref[def:presheaf]{presheaves} of sets \(\mathcal{O} \times \mathcal{F} \to \mathcal{F}\) such that for every open subset \(U \subseteq X\) the map \(\mathcal{O}(U) \times \mathcal{F}(U) \to \mathcal{F}(U)\) defines the structure of an \(\mathcal{O}(U)\)-module on \(\mathcal{F}(U)\).

  A \textbf{morphism} of presheaves of \(\mathcal{O}_{X}\)-modules \(\varphi: \mathcal{F} \to \mathcal{G}\) is a morphism of presheaves of abelian groups \(\varphi: \mathcal{F} \to \mathcal{G}\) such that the following diagram commutes
  \begin{center}
    \begin{tikzcd}
      \mathcal{O}_{X} \times \mathcal{F} \arrow[d, "\id \times \varphi"'] \arrow[r] & \mathcal{F} \arrow[d, "\varphi"] \\
      \mathcal{O}_{X} \times \mathcal{G} \arrow[r] & \mathcal{G}
    \end{tikzcd}
  \end{center}

  We denote the \hyperref[def:category]{category} of presheaves of \(\mathcal{O}_{X}\)-modules over \(X\) as \(\PMod(\mathcal{O}_{X})\).
\end{definition}

\begin{notation}
  When the context is clear, we will drop the \(X\) subindex and say \hyperref[def:presheaf-of-modules]{presheves of \(\mathcal{O}\)-modules} and \(\PMod(\mathcal{O})\).
\end{notation}


\begin{definition} \label{def:restriction-presheaf-of-modules}
  Let \(X\) be a topological space.
  Let \(\mathcal{O}_{1} \to \mathcal{O}_{2}\) be a morphism of \hyperref[def:presheaf]{presheaves} of rings.
  We define the \textbf{restriction} \hyperref[def:functor]{functor}
  \begin{equation*}
    -_{\mathcal{O}_{1}}: \PMod(\mathcal{O}_{2}) \to \PMod(\mathcal{O}_{1})
  \end{equation*}
  by using for any \hyperref[def:presheaf-of-modules]{presheaf of \(\mathcal{O}_{2}\)-modules} \(\mathcal{F}\) the composition map \(\mathcal{O}_{1} \times \mathcal{F} \to \mathcal{O}_{2} \times \mathcal{F} \to \mathcal{F}\).
\end{definition}

\begin{definition} \label{def:tensor-product-presheaf}
  Let \(X\) be a topological space.
  For any \hyperref[def:presheaf]{presheaves} of rings \(\mathcal{O}_{1}, \mathcal{O}_{2}\) define the \textbf{change of rings} \hyperref[def:functor]{functor}
  \begin{equation*}
    \mathcal{O}_{2} \otimes_{\operatorname{pre}, \mathcal{O}_{1}} -: \PMod(\mathcal{O}_{1}) \to \PMod(\mathcal{O}_{2})
  \end{equation*}
  as, for any \hyperref[def:presheaf-of-modules]{presheaf of \(\mathcal{O}_{1}\)-modules} \(\mathcal{G}\) and open subset \(U \subseteq X\)
  \begin{equation*}
    (\mathcal{O}_{2} \otimes_{\operatorname{pre}, \mathcal{O}_{1}} \mathcal{G})(U) \coloneq \mathcal{O}_{2}(U) \otimes_{\mathcal{O}_{1}(U)} \mathcal{G}(U)
  \end{equation*}
\end{definition}

\begin{lemma} \label{lemma:restriction-tensor-adjunction-presheaf}
  For any \(\mathcal{O}_{1} \to  \mathcal{O}_{2}\) morphism of \hyperref[def:presheaf]{presheaves} of rings the \hyperref[def:restriction-presheaf-of-modules]{restriction functor} and the \hyperref[def:tensor-product-presheaf]{change of rings functor} are \hyperref[def:adjoint-functors]{adjoint}. In formula
  \begin{equation*}
    \Hom_{\mathcal{O}_{1}} \left( \mathcal{G}, \mathcal{F}_{\mathcal{O}_{1}} \right) \cong \Hom_{\mathcal{O}_{2}} (\mathcal{O}_{2} \oplus_{\operatorname{pre}, \mathcal{O}_{1}} \mathcal{G}, \mathcal{F})
  \end{equation*}
\end{lemma}
\begin{proof}
  This follows directly from the adjuntion of the corresponding functors in \(\Ring\).
\end{proof}

\begin{definition} \label{def:sheaf-of-modules}
  Let \(X\) be a topological space.
  Let \(\sheaf{O}_{X}\) be a \hyperref[def:sheaf]{sheaf} of rings.
  The \hyperref[def:category]{category} of \textbf{sheaves of \(\sheaf{O}_{X}\)-modules} \(\Mod(\sheaf{O}_{X})\) is the \hyperref[def:subcategory]{full subcategory} of \(\PMod(\iota\sheaf{O}_{X})\) consisting of \hyperref[def:presheaf-of-modules]{presheaves of \(\iota\sheaf{O}_{X}\)-modules} where the the underlying \hyperref[def:presheaf]{presheaves} of abelian modules are \hyperref[def:sheaf]{sheaves}.
\end{definition}

Since, at the end of the day, \hyperref[def:presheaf-of-modules]{presheaves of \(\mathcal{O}\)-modules} and \hyperref[def:sheaf-of-modules]{sheaves of \(\sheaf{O}\)-modules} are \hyperref[def:presheaf]{presheaves} and \hyperref[def:sheaf]{sheaves} of abelian groups respectively, all the constructions we have seen in the past chapters apply.

For example, in the case of \hyperref[def:presheaf-of-modules]{presheaves of \(\mathcal{O}\)-modules} we have.
\begin{lemma}
  Let \(X\) be a topological space and \(x \in X\).
  Let \(\mathcal{O}\) be a \hyperref[def:presheaf]{presheaf} of rings over \(X\).
  Let \(\mathcal{F}\) be a \hyperref[def:presheaf-of-modules]{presheaf of \(\mathcal{O}\)-modules}.

  The canonical map \(\mathcal{O}_{x} \times \mathcal{F}_{x} \to \mathcal{F}_{x}\) coming from the map \(\mathcal{O} \times \mathcal{F} \to \mathcal{F}\) defines an \(\mathcal{O}_{x}\)-module structure on the abelian group \(\mathcal{F}_{x}\).
\end{lemma}

\begin{lemma}
  Let \(X\) be a topological space and \(x \in X\).
  Let \(\mathcal{O} \to \mathcal{O}'\) be a morphism of \hyperref[def:presheaf]{presheaves} of rings over \(X\).
  Let \(\mathcal{F}\) be a \hyperref[def:presheaf-of-modules]{presheaf of \(\mathcal{O}\)-modules} over \(X\).

  A \(\mathcal{O}'_{x}\)-modules we have that
  \begin{align*}
    \mathcal{F}_{x} \otimes_{\mathcal{O}_{x}} \mathcal{O}'_{x} \cong (\mathcal{F} \oplus_{\operatorname{pre}, \mathcal{O}} \mathcal{O}')_{x}
  \end{align*}
\end{lemma}

The case of \hyperref[def:sheafification]{sheafification} is a little bit more tricky.
We have an analogous of \autoref{prop:sheafification-of-algebraic-structure}.

\begin{proposition}
  Let \(X\) be a topological space.
  Let \(\mathcal{O}\) be a \hyperref[def:presheaf]{presheaf} of rings over \(X\) and \(\mathcal{O}^{\#}\) be its \hyperref[def:sheafification]{sheafification}.
  Let \(\mathcal{F}\) be a \hyperref[def:presheaf-of-modules]{presheaf of \(\mathcal{O}\)-modules}.
  Let \(\mathcal{F}^{\#}\) be the \hyperref[def:sheafification]{sheafification} of \(\mathcal{F}\) as a \hyperref[def:presheaf]{presheaf} of abelian groups.

  There exists a map of \hyperref[def:set-valued-sheaf]{sheaves of sets} \(\mathcal{O}^{\#} \times \mathcal{F}^{\#} \to \mathcal{F}^{\#}\) making the diagram
  \begin{center}
    \begin{tikzcd}
      \mathcal{O} \times \mathcal{F} \arrow[r] \arrow[d] & \mathcal{F} \arrow[d] \\
      \mathcal{O}^{\#} \times \mathcal{F}^{\#} \arrow[r] & \mathcal{F}^{\#} 
    \end{tikzcd}
  \end{center}
  commute and which makes \(\mathcal{F}^{\#}\) into a \hyperref[def:sheaf-of-modules]{sheaf of \(\mathcal{O}^{\#}\)-modules}.
  Furthermore, for any \hyperref[def:sheaf-of-modules]{sheaf of \(\mathcal{O}^{\#}\)-modules} \(\sheaf{G}\), all morphisms of \hyperref[def:presheaf-of-modules]{presheaves of \(\mathcal{O}\)-modules} \(\mathcal{F} \to \sheaf{G}\) factor uniquely as \(\mathcal{F} \to \mathcal{F}^{\#} \to \sheaf{G}\) where \(\mathcal{F}^{\#} \to \sheaf{G}\) is a morphism of \hyperref[def:sheaf-of-modules]{sheaves of \(\mathcal{O}^{\#}\)-modules}.

  This defines a sheafification \hyperref[def:functor]{functor}
  \begin{equation*}
    -^{\#}: \PMod(\mathcal{O}) \to \Mod(\mathcal{O}^{\#})
  \end{equation*}
\end{proposition}

\begin{lemma}
  Let \(X\) be a topological space.
  Let \(\mathcal{O}\) be a \hyperref[def:presheaf]{presheaf} of rings over \(X\).

  The \hyperref[def:functor]{functors} \(-^{\#}: \PMod(\mathcal{O}) \to \Mod(\mathcal{O}^{\#})\) and \(\iota: \Mod(\sheaf{O}) \to \PMod(\iota\sheaf{O})\) are \hyperref[def:adjoint-functors]{adjoint}.
  In formula
  \begin{equation*}
    \Hom_{\PMod(\mathcal{O})} (\mathcal{F}, \iota\sheaf{G}) \cong \Hom_{\Mod(\mathcal{O}^{\#})}(\mathcal{F}^{\#}, \sheaf{G})
  \end{equation*}
\end{lemma}

Using this adjuntion we can \textit{translate} \autoref{def:restriction-presheaf-of-modules} and \autoref{def:tensor-product-presheaf}.
\begin{definition}
  Let \(X\) be a topological space.
  Let \(\sheaf{O}_{1} \to \sheaf{O}_{2}\) be a morphism of \hyperref[def:sheaf]{sheaves} of rings.
  Let \(\sheaf{F}\) be a \hyperref[def:sheaf-of-modules]{sheaf of \(\sheaf{O}_{2}\)-modules}.

  We define the \textbf{restriction} of \(\sheaf{F}\) to \(\sheaf{O}_{1}\) as
  \begin{align*}
    \sheaf{F}_{\sheaf{O}_{1}} \coloneq \left( \iota \left( \sheaf{F}_{\sheaf{O}_{1}} \right) \right)^{\#} = \sheaf{F}_{\sheaf{O}_{1}}
  \end{align*}
  where \(\iota(\sheaf{F}_{\sheaf{O}_{1}})\) is the \hyperref[def:restriction-presheaf-of-modules]{restriction as a presheaf of \(\iota\sheaf{O}_{2}\)-modules}.

  I.e: the restriction of \hyperref[def:sheaf-of-modules]{sheaves of modules} coincides with the \hyperref[def:restriction-presheaf-of-modules]{restriction of presheaves of modules}.
  This induces a \hyperref[def:functor]{functor}
  \begin{equation*}
    -_{\sheaf{O}_{1}}: \Mod(\sheaf{O}_{2}) \to \Mod(\sheaf{O}_{1})
  \end{equation*}
\end{definition}

\begin{definition}
  Let \(X\) be a topological space.
  Let \(\sheaf{O}_{1}, \sheaf{O}_{2}\) be \hyperref[def:sheaf]{sheaves} of rings over \(X\).
  Let \(\sheaf{F}\) be a \hyperref[def:sheaf-of-modules]{sheaf of \(\sheaf{O}_{2}\)-modules}.

  We define the \textbf{change of basis} of \(\sheaf{F}\) into \(\sheaf{O}_{2}\) as
  \begin{equation*}
    \sheaf{O}_{2} \otimes_{\sheaf{O}_{1}} \sheaf{F} \coloneq \left( \iota\sheaf{O}_{2} \otimes_{\operatorname{pre}, \iota\sheaf{O}_{1}} \iota\sheaf{F} \right)^{\#}
  \end{equation*}
  This induces a \hyperref[def:functor]{functor}
  \begin{equation*}
    \sheaf{O}_{2} \otimes_{\sheaf{O}_{1}} -: \Mod(\sheaf{O}_{1}) \to \Mod(\sheaf{O}_{2})
  \end{equation*}
\end{definition}

We also have an analogous of \autoref{lemma:restriction-tensor-adjunction-presheaf}
\begin{lemma} \label{lemma:restriction-tensor-adjunction-sheaf}
  Let \(X\) be a topological space.
  Let \(\sheaf{O}_{1} \to \sheaf{O}_{2}\) be a morphism of \hyperref[def:sheaf]{sheaves} of rings over \(X\).

  The \hyperref[def:functor]{functors} \(-_{\sheaf{O}_{1}}\) and \(\sheaf{O}_{2} \otimes_{\sheaf{O}_{1}} -\) are \hyperref[def:adjoint-functors]{adjoint}.
  In formula, if \(\sheaf{F}\) is a \hyperref[def:sheaf-of-modules]{sheaf of \(\sheaf{O}_{2}\)-modules} and \(\sheaf{G}\) is a \hyperref[def:sheaf-of-modules]{sheaf of \(\sheaf{O}_{1}\)-modules} we have
  \begin{equation*}
    \Hom_{\Mod(\sheaf{O}_{1})} \left( \sheaf{G}, \sheaf{F}_{\sheaf{O}_{1}} \right) \cong \Hom_{\Mod(\sheaf{O}_{2})} \left( \sheaf{O}_{2} \otimes_{\sheaf{O}_{1}} \sheaf{G}, \sheaf{F} \right)
  \end{equation*}
\end{lemma}

\begin{lemma}
  Let \(X\) be a topological space and \(x \in X\).
  Let \(\sheaf{O} \to \sheaf{O}'\) be a morphism of \hyperref[def:sheaf]{sheaves} of rings over \(X\).
  Let \(\sheaf{F}\) be a \hyperref[def:sheaf-of-modules]{sheaf of \(\sheaf{O}\)-modules}.

  As \(\sheaf{O}'_{x}\)-modules we have
  \begin{align*}
    \sheaf{F}_{x} \otimes_{\sheaf{O}_{x}} \sheaf{O}'_{x} \cong \left( \sheaf{F} \otimes_{\sheaf{O}} \sheaf{O}' \right)_{x}
  \end{align*}
\end{lemma}

\begin{theorem} \label{thm:mod-ox-is-abelian}
  \(\Mod(\sheaf{O})_{X}\) is an \hyperref[def:abelian-category]{abelian category}
\end{theorem}
\begin{proof}
  In \autoref{prop:limits-colimits-presheaf} we have seen that \hyperref[def:categorical-limit]{limits} and \hyperref[def:categorical-colimit]{colimits} of \hyperref[def:presheaf]{presheaves} can be computed \textit{open set by open set}.
  Therefore since the \hyperref[def:category]{category} of modules over a ring is \hyperref[def:abelian-category]{abelian}, the category \(\PMod(\sheaf{O}_{X})\) inherits this property.

  Let's see how we can see that \(\Mod(\sheaf{O}_{X})\) is also \hyperref[def:abelian-category]{abelian}.
  It is clear that \(\Mod(\sheaf{O}_{X})\) is an \hyperref[def:additive-category]{additive category}.
  Therefore, we only need to check that \hyperref[def:categorical-kernel]{kernels} and \hyperref[def:categorical-cokernel]{cokernels} exist in \(\Mod(\sheaf{O}_{X})\).

  Let \(\varphi: \sheaf{F} \to \sheaf{G}\) be a morphism of \hyperref[def:sheaf-of-modules]{sheaves of \(\sheaf{O}\)-modules}.
  We have an associated morphism of \hyperref[def:presheaf-of-modules]{presheaves of \(\iota\sheaf{O}\)-modules} \(\iota\varphi: \iota\sheaf{F} \to \iota\sheaf{G}\).
  Using that \(\PMod(\iota\sheaf{O}_{X})\) is an \hyperref[def:abelian-category]{abelian category} the morphism \(\iota\varphi\) admits a \hyperref[def:categorical-kernel]{kernel} \(\ker_{\operatorname{pre}} \iota \varphi\) and \hyperref[def:categorical-cokernel]{cokernel} \(\coker_{\operatorname{pre}} \iota\varphi\) in \(\PMod(\iota\sheaf{O}_{X})\).
  Therefore we can define the \hyperref[def:categorical-kernel]{kernel} and \hyperref[def:categorical-cokernel]{cokernel} of \(\varphi\) using the \hyperref[def:sheafification]{sheafification functor}
  \begin{align*}
    \ker \varphi &\coloneq \left( \ker_{\operatorname{pre}} \iota\varphi \right)^{\#} = \ker_{\operatorname{pre}} \left( \iota \varphi \right)^{\#} = \ker_{\operatorname{pre}} \varphi \\
    \coker \varphi &\coloneq \left( \operatorname{coker}_{\operatorname{pre}} \iota\varphi \right)^{\#}
  \end{align*}
  where the equalities in \(\ker \varphi\) come from the fact that \hyperref[prop:adjoint-properties]{right adjoints preserve limits}.
\end{proof}
\end{document}
